\chapter{MAPPING: LIDAR SLAM}
\label{ch:slam}

The map is built on the laptop from the board's \topic{/scan} and the
\frm{odom}$\to$\frm{base\_link} transform. Two SLAM packages are set up:
gmapping, a Rao--Blackwellized particle filter that uses odometry and is the
default, and hector\_mapping, a scan matcher that can run on the lidar alone.

\section{OCCUPANCY GRID}

Both build a 2-D occupancy grid with \SI{5}{cm} cells. Each cell $m_i$ holds
the log-odds of being occupied,
\begin{equation}
\ell_i = \log\frac{p(m_i)}{1-p(m_i)},
\end{equation}
and every scan updates the cells along each beam by the inverse sensor model
\cite{thrun2005probabilistic}:
\begin{equation}\label{eq:logodds}
\ell_{i,t} = \ell_{i,t-1} + \log\frac{p(m_i\mid z_t,x_t)}{1-p(m_i\mid z_t,x_t)} - \ell_0 .
\end{equation}
Cells a beam passes through become more likely free, and the cell where it
ends more likely occupied. Because the update is additive, a scan registered
at the wrong pose does not erase the old walls; it adds a second, rotated or
shifted copy. That is why a pose error shows up as a ``jumbled'' map rather
than a moved one.

\section{GMAPPING}
\label{sec:gmapping}

\subsection{Rao--Blackwellized particle filter}

SLAM estimates the joint posterior of the robot's path $x_{1:t}$ and the map
$m$, given scans $z_{1:t}$ and odometry $u_{1:t-1}$. gmapping
\cite{grisetti2007improved} factors it as
\begin{equation}\label{eq:rbpf}
p(x_{1:t},m\mid z_{1:t},u_{1:t-1}) =
p(m\mid x_{1:t},z_{1:t})\;p(x_{1:t}\mid z_{1:t},u_{1:t-1}).
\end{equation}
Given a path, the map follows by \eqref{eq:logodds}, so only the path needs a
particle filter; each of the $M$ particles carries one path hypothesis and
its own map \cite{doucet2000rao,murphy1999bayesian}. Per accepted scan, each
particle $i$:
\begin{enumerate}
\item \textbf{Predicts} with the odometry motion model: the change in
  \frm{odom}$\to$\frm{base\_link} since the last update, split into a
  translation $\Delta_\rho$ and a rotation $\Delta_\psi$, is perturbed by
  Gaussian noise with standard deviations
\begin{equation}\label{eq:gmapmotion}
\sigma_\rho = s_{rr}\abs{\Delta_\rho} + s_{rt}\abs{\Delta_\psi},\qquad
\sigma_\psi = s_{tr}\abs{\Delta_\rho} + s_{tt}\abs{\Delta_\psi},
\end{equation}
  (parameters \texttt{srr}, \texttt{srt}, \texttt{str}, \texttt{stt}).
\item \textbf{Scan-matches} the new scan against its own map, starting from
  the predicted pose $x'$, by hill-climbing in steps of \texttt{lstep} (m) and
  \texttt{astep} (rad) to maximise a likelihood-field score
\begin{equation}\label{eq:gmapscore}
s(x) = \sum_{j}\exp\!\left(-\frac{d_j(x)^2}{\sigma}\right),
\end{equation}
  where $d_j(x)$ is the distance from the endpoint of beam $j$ (projected from
  pose $x$) to the nearest occupied cell, and $\sigma$ is the parameter
  \texttt{sigma}. Only beams shorter than \texttt{maxUrange} are used.
\item \textbf{Samples} its new pose from a Gaussian fitted around the best
  match $\hat x^\star$ (the ``improved proposal''), which combines the scan
  match with \eqref{eq:gmapmotion}. If the best score is below
  \texttt{minimumScore}, the match is considered unreliable and the
  particle falls back to the odometry prediction $x'$.
\item \textbf{Reweights}: $w^{(i)}_t \propto w^{(i)}_{t-1}\,\eta^{(i)}$, where
  $\eta^{(i)}=\int p(z_t\mid m^{(i)},x)\,p(x\mid x^{(i)}_{t-1},u_{t-1})\,\dd x$
  is approximated from the same samples.
\end{enumerate}
When the effective sample size
\begin{equation}
N_{\mathrm{eff}} = \frac{1}{\sum_i \bigl(\tilde w^{(i)}\bigr)^2}
\end{equation}
of the normalised weights $\tilde w$ falls below
$\texttt{resampleThreshold}\cdot M$, particles are resampled. The best
particle's map is published on \topic{/map} every
\texttt{map\_update\_interval} seconds, and its pose sets the transform
\frm{map}$\to$\frm{odom} such that
\begin{equation}
T_{\frm{map}\leftarrow\frm{base\_link}} =
T_{\frm{map}\leftarrow\frm{odom}}\;T_{\frm{odom}\leftarrow\frm{base\_link}} .
\end{equation}
That composition is how the SLAM pose and the odometry pose integrate: the
odometry provides a smooth, \SI{30}{Hz} pose that drifts, and SLAM publishes
the slowly changing correction that cancels the drift.

\subsection{When scans are processed}

A scan is processed only if the robot has moved at least
\texttt{linearUpdate} metres or turned \texttt{angularUpdate} radians
according to odometry since the last processed scan, or, if
\texttt{temporalUpdate} is positive, if that many seconds have passed.

\subsection{Parameters}

\begin{table}[htbp]
\centering
\caption{gmapping parameters (\file{laptop/bbai\_slam/gmapping\_bbai.launch})}
\label{tab:gmap}
\small
\begin{tabular}{lccl}
\toprule
Parameter & 1 Oct & 2 Oct & Effect \\
\midrule
\texttt{particles} $M$ & 50 & 40 & path hypotheses \\
\texttt{temporalUpdate} & \SI{2.0}{s} & off ($-1$) & process scans while parked \\
\texttt{linearUpdate} / \texttt{angularUpdate} & 0.10 / 0.10 & 0.10 / 0.10 & motion between scans (m / rad) \\
\texttt{minimumScore} & 0 & 100 & reject weak scan matches \\
\texttt{srr}, \texttt{srt}, \texttt{str}, \texttt{stt} & 0.1, 0.2, 0.1, 0.2$^\dagger$ & 0.05 each & odometry noise \eqref{eq:gmapmotion} \\
\texttt{sigma}, \texttt{lstep}, \texttt{astep} & 0.05$^\dagger$ & 0.05 & scan-match kernel and steps \\
\texttt{maxUrange} / \texttt{maxRange} & 8 / 12 m & 6 / 12 m & usable / maximum beam length \\
\texttt{map\_update\_interval} & \SI{2.0}{s} & \SI{1.0}{s} & map publishing \\
\texttt{resampleThreshold} & 0.5$^\dagger$ & 0.5 & resample when $N_{\mathrm{eff}}<M/2$ \\
\texttt{delta} & \SI{0.05}{m} & \SI{0.05}{m} & cell size \\
\texttt{transform\_publish\_period} & \SI{0.05}{s}$^\dagger$ & \SI{0.05}{s} & \frm{map}$\to$\frm{odom} rate \\
\bottomrule
\multicolumn{4}{l}{\footnotesize $^\dagger$ package default, not set in the launch file}
\end{tabular}
\end{table}

\influx{The 2 October values are on branch \texttt{slam-estimator-updates}
and were being tested when this was written.}

\section{HECTOR MAPPING}
\label{sec:hector}

hector\_mapping \cite{kohlbrecher2011flexible} has no particles and no
odometry model. It aligns each scan directly to the map built so far. With the
scan's beam endpoints $\bm s_j$ in the robot frame and the pose
$\bm\xi=(p_x,p_y,\psi)$, the endpoints in the map are
\begin{equation}
\bm S_j(\bm\xi) = R(\psi)\bm s_j + \begin{bmatrix}p_x\\p_y\end{bmatrix},
\end{equation}
and the pose minimises
\begin{equation}\label{eq:hectorobj}
\bm\xi^\star=\argmin_{\bm\xi}\sum_j\bigl[1-M(\bm S_j(\bm\xi))\bigr]^2,
\end{equation}
where $M(\cdot)\in[0,1]$ is the occupancy value bilinearly interpolated
between grid cells. One Gauss--Newton step from the current estimate is
\begin{equation}\label{eq:hectorgn}
\Delta\bm\xi = H^{-1}\sum_j\left[\nabla M(\bm S_j)\frac{\partial\bm S_j}{\partial\bm\xi}\right]\T
\bigl[1-M(\bm S_j)\bigr],\qquad
H=\sum_j\left[\nabla M\frac{\partial\bm S_j}{\partial\bm\xi}\right]\T
\left[\nabla M\frac{\partial\bm S_j}{\partial\bm\xi}\right],
\end{equation}
with
\begin{equation}
\frac{\partial\bm S_j}{\partial\bm\xi}=
\begin{bmatrix}1&0&-\sin\psi\,s_{j,x}-\cos\psi\,s_{j,y}\\
0&1&\cos\psi\,s_{j,x}-\sin\psi\,s_{j,y}\end{bmatrix}.
\end{equation}
The match is run coarse-to-fine on a pyramid of maps (2 levels here) to widen
its basin of convergence. Hector converges only if the true pose is within
about one coarse cell and a few degrees of the previous estimate, so it fails
when the robot turns faster than the scan rate can follow
(\S\ref{sec:skew}). On this robot it is the fallback for mapping without
odometry (\file{lidar\_only\_bbai.launch}), for example when carrying the
robot by hand.

\section{SCAN RATE AND ROTATION}
\label{sec:skew}

The A1 delivers about 7.2 scans per second, so between scans, and during the
\SI{0.14}{s} it takes to sweep one scan, the robot turns by
\begin{equation}
\Delta\psi_{\mathrm{scan}} = \frac{\omega}{f_{\mathrm{scan}}}.
\end{equation}
\begin{center}
\small
\begin{tabular}{lccc}
\toprule
Turn rate & \SI{0.5}{rad/s} & \SI{1.5}{rad/s} & \SI{6}{rad/s} \\
\midrule
Rotation per scan & \SI{4}{\degree} & \SI{12}{\degree} & \SI{48}{\degree} \\
\bottomrule
\end{tabular}
\end{center}
Neither SLAM package corrects the distortion within a scan (each beam is
treated as if taken at the scan's time stamp), so a scan taken while turning
at \SI{1.5}{rad/s} is itself bent by up to \SI{12}{\degree} from its first
beam to its last. gmapping's \texttt{angularUpdate} of \SI{0.1}{rad}
(\SI{5.7}{\degree}) also means that at \SI{1.5}{rad/s} nearly every scan is
processed while bent. Mapping is most reliable with turns below about
\SI{0.5}{rad/s}. The new teleop curve keeps half stick at
\SI{1.5}{rad/s}; turning gently (well under half stick) while mapping is
better still.

\section{DIAGNOSING THE SLAM DRIFT OF 1 OCTOBER}
\label{sec:slamdiag}

Jason observed that gmapping made a good first map, lost track once the robot
moved, and, when left stationary, slowly rotated the map into a jumble. Each
symptom has a distinct likely cause.

\paragraph{Map rotating while parked.} With \texttt{temporalUpdate} at
\SI{2}{s}, gmapping processed a scan every two seconds even when the robot did
not move. The gyroscope heading drifted about \SI{1.6}{\degree/min} while
parked (\S\ref{sec:gyro}), so each new scan arrived with a slightly rotated
odometry pose. The scan matcher absorbs a small rotation, but with only
about 360 beams per scan and a cell size of \SI{5}{cm}, its correction is
imperfect, and the residual is written into the map by \eqref{eq:logodds}.
Repeated every two seconds this smears the walls into a rotating blur. Two
changes address it (branch \texttt{slam-estimator-updates}):
\texttt{temporalUpdate} $=-1$, so a parked robot adds no scans, and the
zero-velocity update (\S\ref{sec:zupt}), which removes the heading drift
while parked. Parked drift measured afterwards was zero.

\paragraph{Losing track while driving.} Several causes can each produce this,
and they are best ruled out in order:
\begin{enumerate}
\item \textbf{Lidar mounting yaw (found and fixed).} The static transform
  \frm{base\_link}$\to$\frm{laser} had zero yaw. \texttt{rplidar\_ros}
  reports angle zero along the A1's own zero direction, which on this robot
  points backward. With the yaw off by \SI{180}{\degree}, a spin in place
  still looks consistent (the scan rotates in the same direction either
  way), but driving forward makes the walls appear to approach from behind:
  odometry and scan disagree on every metre travelled and the map breaks as
  soon as the robot drives, which is exactly the symptom.
  \texttt{lidar\_mount\_test.py} confirmed it on 2 October 2026: during a
  turn the scan shifted by $-\Delta\psi$ (a rotation, not a mirror image),
  and driving forward shortened the ranges most at bearing \SI{180}{\degree}.
  The transform now has yaw $\pi$ in \file{robot.launch} and in the laptop's
  \file{lidar\_only\_bbai.launch}. Since the fix, the \frm{map}$\to$\frm{odom}
  correction stays within a few degrees and centimetres across
  $\pm\SI{90}{\degree}$ turns and \SI{0.3}{m} drives. The lidar's position
  is still assumed to be at the robot's centre until it is measured.
\item \textbf{Turning too fast.} See \S\ref{sec:skew}.
\item \textbf{Weak scan matches.} At the A1's range and with about 360 beams,
  narrow corridors and open rooms give the matcher little to hold on to.
  \texttt{minimumScore} now rejects weak matches in favour of odometry, which
  is reliable over short distances after the calibration
  (\SI{0.914}{m} measured for \SI{0.914}{m}). \texttt{maxUrange} was reduced to
  \SI{6}{m} because the A1's error grows with range.
\item \textbf{Odometry noise model.} With \texttt{srr}--\texttt{stt} at their
  defaults (0.1--0.2), particles spread far more than the calibrated odometry
  warrants, and with only 40--50 particles the filter can lock onto a wrong
  hypothesis. They are now 0.05.
\item \textbf{Clocks.} Scans and transforms are stamped by the board's clock,
  which has no battery-backed time source and is not synchronised to the
  laptop. gmapping looks up odometry at each scan's stamp, which both come
  from the board, so this is not expected to cause the drift, but running
  \texttt{chrony} on both machines removes a class of TF extrapolation
  warnings.
\end{enumerate}

\influx{The yaw fix, the mount test and a \file{reset\_map.sh} that waits for
the old gmapping node to exit before relaunching are on branch
\texttt{slam-estimator-updates} (commit 0044110), not yet on \texttt{main}.
Longer mapping runs have not been tested yet.}
